\documentclass[11pt]{amsart}

\usepackage{amssymb}
\usepackage{mathtools}
\usepackage[margin=1.1in]{geometry}
\usepackage[hidelinks]{hyperref}

\newtheorem{theorem}{Theorem}
\newtheorem*{theorem*}{Theorem}
\newtheorem{lemma}[theorem]{Lemma}
\newtheorem{proposition}[theorem]{Proposition}
\theoremstyle{remark}
\newtheorem*{remark}{Remark}

\newcommand{\Z}{\mathbb{Z}}
\newcommand{\E}{\mathbb{E}}
\DeclarePairedDelimiter{\abs}{\lvert}{\rvert}
\DeclarePairedDelimiter{\floor}{\lfloor}{\rfloor}
\DeclarePairedDelimiter{\ceil}{\lceil}{\rceil}

\title[Theorem 8 from a weaker form of Theorem 5]{Large sumsets from medium-sized subsets:\\
Theorem 8 from a weaker form of Theorem 5}
\author{Gabriel Dahia}
\date{}

\begin{document}

\begin{abstract}
Bollobás, Leader and Tiba prove Theorem 1 of \emph{Large sumsets from medium-sized subsets} from
two results that they quote: Theorem 4, a stability form of the Cauchy--Davenport theorem, and
Theorem 5, Theorem 31 of their paper \emph{Large sumsets from small subsets}. Theorem 5 gives a
family of bounded subsets of a set $B$ that works on average for every set $A$ close to an
interval. We show that the proof of Theorem 8, and hence of Theorem 1, goes through with a weaker
statement, in which a single bounded subset of $B$ is chosen for each set $A$. Only the sample of
$Y$ in the proof of Theorem 8 changes, and the constant grows by a bounded factor. This is the
route followed by the Lean formalization.
\end{abstract}

\maketitle

\section{Notation}

For finite sets $A$ and $B$ of an abelian group, $A + B = \{a + b : a \in A,\ b \in B\}$ and
$\abs{A}$ is the number of elements of $A$. We write $A \mathbin{\triangle} B$ for the symmetric
difference, $\Z_p$ for the integers modulo a prime $p$, and $[a, b)$ for the set of integers $u$
with $a \le u < b$. An interval of $\Z_p$ with $m$ elements is a set
$\{x, x + 1, \dots, x + m - 1\}$. Theorem numbers refer to \emph{Large sumsets from
medium-sized subsets} unless stated otherwise; ``the small paper'' is \emph{Large sumsets from
small subsets}.

For $d \ge 1$, a residue $r \in \Z_d$ and a finite set $S \subseteq \Z$, the class of $S$ at $r$ is
$S_r = \{s \in S : s \equiv r \pmod d\}$. For sets $S, T \subseteq \Z$ and $z \in \Z_d$,
\begin{equation}\label{eq:classes}
  (S + T)_z = \bigcup_{x + y = z} (S_x + T_y),
  \qquad
  \abs{S + T} = \sum_{z \in \Z_d} \abs{(S + T)_z}.
\end{equation}
If $\bar x$ and $\bar y$ are fixed representatives of $x$ and $y$, the maps
$s \mapsto (s - \bar x)/d$ on $S_x$ and $s \mapsto (s - \bar y)/d$ on $T_y$ are injective and
$\abs{S_x + T_y}$ is the size of the sumset of the images. The class at $r$ of an interval of
length $m$ is mapped onto an interval of length between $\floor{m/d}$ and $\floor{m/d} + 1$, so
the classes of a set close to an interval are again sets of integers close to intervals.

\section{Theorems 4 and 5}

Theorem 4 of the paper is the following statement. It is Corollary 24 of the small paper, which
the formalization of the small paper proves.

\begin{theorem*}[Theorem 4]
For $\beta, \gamma > 0$ there is $\varepsilon > 0$ with the following property. Let
$A, B \subseteq \Z_p$ have at least two elements, with $\abs{A} + \abs{B} \le (1 - \beta) p$ and
$\abs{A + B} \le \abs{A} + \abs{B} - 1 + \varepsilon \min(\abs{A}, \abs{B})$. Then $A$ and $B$
are contained in arithmetic progressions $P$ and $Q$ with the same common difference such that
$\abs{P \setminus A}, \abs{Q \setminus B} \le \gamma \min(\abs{A}, \abs{B})$.
\end{theorem*}

Theorems 1 and 6 use Theorem 4 exactly as stated, so nothing changes there. Theorem 5 is Theorem
31 of the small paper.

\begin{theorem*}[Theorem 5]
For $\beta > 0$ there is $\gamma > 0$ such that for $\alpha > 0$ there is $c$ with the following
property. Let $I$ and $J = \{y, \dots, y + n - 1\}$ be intervals of $\Z_p$ with $m$ and $n$
elements, $\alpha n \le m \le \alpha^{-1} n$ and $m + n \le (1 - \beta) p$, and let
$B \subseteq J$ contain $y$ and $y + n - 1$ and miss at most $\gamma \min(m, n)$ points of $J$.
Then there is a family $\mathcal F$ of subsets of $B$ with at most $c$ elements each such that
every $A \subseteq \Z_p$ with $\abs{A \mathbin{\triangle} I} \le \gamma \min(m, n)$ satisfies
\[
  \E_{F \in \mathcal F} \abs{A + F} \ge \abs{A} + \abs{B} - 1 .
\]
\end{theorem*}

The point of Theorem 5 is that $\mathcal F$ depends on $B$ but not on $A$. The formalization of
the small paper proves the following weaker statement, which is the form used there for Theorem
27 of the small paper. We state it for intervals of integers, the setting of the proof of
Theorem 8.

\begin{theorem}[weak form of Theorem 5]\label{thm:weak5}
Let $\gamma = 1/1600$. For $\alpha > 0$ there is $c$ with the following property. Let
$I = [x, x + m)$ and $J = [y, y + n)$ with $\alpha n \le m \le \alpha^{-1} n$, let $B \subseteq J$
contain $y$ and $y + n - 1$ with $\abs{J \setminus B} \le \gamma \min(m, n)$, and let
$A \subseteq \Z$ be finite with $\abs{A \mathbin{\triangle} I} \le \gamma \min(m, n)$. Then there
is $F \subseteq B$ with $\abs{F} \le c$ and
\[
  \abs{A + F} \ge \abs{A} + \abs{B} - 1 .
\]
\end{theorem}

Theorem 5 implies the analogous statement in $\Z_p$, since some member of $\mathcal F$ is at least
the average. The difference is the order of the quantifiers: in Theorem~\ref{thm:weak5} the set
$F$ depends on $A$.

\begin{proof}[Proof of Theorem~\ref{thm:weak5} from the small paper]
Put $\mu = \min(m, n)$. Since $\abs{A \mathbin{\triangle} I}$ and $\abs{J \setminus B}$ are at most
$\gamma \mu$, we have $m \le \abs{A} + \gamma\mu$, $n \le \abs{B} + \gamma\mu$,
$\mu / 2 \le \min(\abs{A}, \abs{B})$ and therefore
$\gamma \mu \le \delta \min(\abs{A}, \abs{B})$ with $\delta = 1/800$. Let
$\alpha' = \min(\alpha, 1)/2$, so that $\alpha' \abs{B} \le \min(\abs{A}, \abs{B})$, and let
$c = \ceil{64/\alpha'}^2 + 2\ceil{64/\alpha'} + 10$.

If $\abs{B} \le c$, take $F = B$: the Cauchy--Davenport inequality in $\Z$ gives
$\abs{A + B} \ge \abs{A} + \abs{B} - 1$. Otherwise $\abs{A}, \abs{B} \ge 2$. Choose $R$ with
$A \subseteq [x - R, x + m + R)$ and a prime $q > 2(\abs{A} + \abs{B} + m + n + 2R) + 1$.
Reduction modulo $q$ is injective on $[x - R, x + m + R)$, on $J$ and on
$[x + y - R, x + y + m + n + R) \supseteq A + J$, so it preserves all the sizes above, maps $I$
and $J$ to intervals of $\Z_q$, and gives $\abs{A} + \abs{B} \le q/2$. The form of Theorem 31
proved in the small paper, applied in $\Z_q$ with $\beta = 1/2$, $\delta = 1/800$ and $\alpha'$,
gives $F \subseteq B$ with $\abs{F} \le c$ and $\abs{A} + n \le \abs{A + F} + 1$. As
$n \ge \abs{B}$, this is the claim.
\end{proof}

\section{Theorem 8}

\begin{theorem*}[Theorem 8]
There is $\gamma > 0$ such that for $\alpha > 0$ there is $c$ with the following property. Let
$X \subseteq I$ and $Y \subseteq J$, where $I$ and $J$ are intervals of $\Z$ (or of $\Z_p$, with
$m + n \le p$) with $m$ and $n$ elements, $\alpha n \le m \le \alpha^{-1} n$ and
$\abs{I \setminus X}, \abs{J \setminus Y} \le \gamma \min(m, n)$. If
$1 \le c_1 \le \abs{X}$ and $1 \le c_2 \le \abs{Y}$ satisfy
$c_1 c_2 \ge c \max(\abs{X}, \abs{Y})$, there are $X' \subseteq X$ and $Y' \subseteq Y$ with
$\abs{X'} = c_1$, $\abs{Y'} = c_2$ and $\abs{X' + Y'} \ge \abs{X} + \abs{Y} - 1$.
\end{theorem*}

The case of $\Z_p$ follows from the case of $\Z$ by lifting the intervals, since $m + n \le p$. We
recall the proof of the paper in $\Z$ and indicate the one change.

\subsection*{Setup}
We may assume $\alpha \le 1$ and $c_2 \le c_1$; the other case is symmetric. Let $\gamma_0 \le 1$
and $c_0$ be the constants of Theorem~\ref{thm:weak5} for $\alpha/2$, choose $t$ with
$2^{-t} \le \alpha\gamma_0/128$, and put $\gamma = \gamma_0/256$. Choose $d \ge 1$ of order
$c_2/(K + 1)$, where $K$ is a constant fixed below, so that
\[
  \theta = \frac{\gamma_0 \min(m, n)}{d} \ge 128 .
\]
Call a class $X_r$ \emph{dense at level $\lambda$} if it misses at most $\lambda$ points of the
corresponding class of $I$, and similarly for $Y$. Let $E_X$ and $E_Y'$ be the sets of residues
whose classes of $X$ and of $Y$ are dense at level $\theta/16$, and let $E \supseteq E_Y'$ be the
set of residues whose classes of $Y$ are dense at level $\theta/4$. Since $X$ misses at most
$\gamma \min(m, n) = d\theta/256$ points of $I$, Markov's inequality gives
$\abs{\Z_d \setminus E_X} \le d/16$, and likewise $\abs{\Z_d \setminus E_Y'} \le d/16$. Let
$r_0$ be a residue with $\abs{X_{r_0}}$ largest; then $r_0 \in E_X$. Every class has at most
$m/d + 2$ or $n/d + 2$ elements, and a dense class at least $m/d - 2 - \theta/16$ or
$n/d - 2 - \theta/16$. Hence, for $x \in E_X$, $y \in E_Y'$ and $y' \notin E$,
\begin{equation}\label{eq:gain}
  \abs{X_x} + \abs{Y_y} \ge \abs{X_{r_0}} + \abs{Y_{y'}} + g,
  \qquad g = \theta/16 ,
\end{equation}
because $\theta \ge 128$.

\subsection*{The subsets given by Theorem 5}
Let $x \in E_X$ and $y \in E$. After the rescaling of Section~1, $Y_y$ lies in an interval of
length about $n/d$ whose endpoints it contains and of which it misses at most $\theta/4$ points,
and $X_x$ differs from an interval of length about $m/d$ in at most $\theta/16 + O(1)$ points.
(The paper's Theorem 5 needs the endpoints of $J$ in $B$; a dense class is first shrunk to the
interval between its least and largest elements, which changes the length by at most $\theta/4$.)
The lengths satisfy $\tfrac{\alpha}{2}\nu \le \mu \le (\tfrac{\alpha}{2})^{-1}\nu$ and
$\theta/4 \le \gamma_0 \min(\mu, \nu)$, where $\mu$ and $\nu$ are the two lengths. So
Theorem~\ref{thm:weak5} gives a set
\[
  G(x, y) \subseteq Y_y, \qquad \abs{G(x, y)} \le c_0, \qquad
  \abs{X_x + G(x, y)} \ge \abs{X_x} + \abs{Y_y} - 1 .
\]
For the other pairs put $G(x, y) = \emptyset$.

The paper uses Theorem 5 here to get, for each $y \in E$, a family $\mathcal F_y$ of subsets of
$Y_y$ that works on average for every $x \in E_X$ at once, and samples one member of each
$\mathcal F_y$, independently of the sample of $X$. Theorem~\ref{thm:weak5} gives no such family.
Instead, the sample of $Y$ below is built from the sample of $X$.

\subsection*{The samples}
Let $\xi = (\xi_1, \dots, \xi_t)$ be uniform in $E_X^t$ and put
$S(\xi) = \{r_0, \xi_1, \dots, \xi_t\} \subseteq \Z_d$. The samples are
\[
  X_\xi = \bigcup_{x \in S(\xi)} X_x ,
  \qquad
  Y_\xi = \bigcup_{y \in E}\ \bigcup_{x \in S(\xi)} G(x, y) .
\]
The paper's sample of $X$ is the same. In the paper, the part of the sample of $Y$ in the class
$y$ is one random member of $\mathcal F_y$; here it is the union of the sets $G(x, y)$ over the
$t + 1$ classes $x$ of the sample of $X$. Thus
\[
  \abs{X_\xi} \le (t + 1)\Bigl(\frac{m}{d} + 2\Bigr),
  \qquad
  \abs{Y_\xi} \le d\,(t + 1)\,c_0 ,
\]
where the paper has $d c_0$ in the second bound.

\begin{lemma}[Lemma F\&G, adapted]\label{lem:fg}
$\E_\xi \abs{X_\xi + Y_\xi} \ge \abs{X} + \abs{Y} - d$.
\end{lemma}

\begin{proof}
By \eqref{eq:classes}, $\abs{X_\xi + Y_\xi} = \sum_z \abs{(X_\xi + Y_\xi)_z}$. Fix $z \in \Z_d$ and
put $s_z = \abs{X_{r_0}} + \abs{Y_{z - r_0}}$. We show
\begin{equation}\label{eq:class}
  \E_\xi \abs{(X_\xi + Y_\xi)_z} \ge s_z - 1 .
\end{equation}
The only property of $Y_\xi$ used is that, for $x \in S(\xi)$ and $y \in E$ with $x + y = z$,
\begin{equation}\label{eq:pair}
  \abs{(X_\xi + Y_\xi)_z} \ge \abs{X_x + G(x, y)} ,
\end{equation}
which holds by \eqref{eq:classes} because $X_x \subseteq X_\xi$ and $G(x, y) \subseteq Y_\xi$.
This is where the sample of $Y$ must contain $G(x, y)$ for the classes $x$ of the sample of $X$.

If $z - r_0 \in E$, then \eqref{eq:pair} with $x = r_0 \in S(\xi)$ and $y = z - r_0$ gives
$\abs{(X_\xi + Y_\xi)_z} \ge s_z - 1$ for every $\xi$.

If $z - r_0 \notin E$, call $\xi$ good when $z - \xi_i \in E_Y'$ for some $i$. For good $\xi$,
\eqref{eq:pair} with $x = \xi_i$ and $y = z - \xi_i$, together with \eqref{eq:gain} for
$y' = z - r_0$, gives
\[
  \abs{(X_\xi + Y_\xi)_z} \ge \abs{X_{\xi_i}} + \abs{Y_{z - \xi_i}} - 1 \ge s_z + g - 1 .
\]
The residues $x \in E_X$ with $z - x \notin E_Y'$ number at most
$\abs{\Z_d \setminus E_Y'} \le d/16 \le \abs{E_X}/2$, so $\xi$ is bad with probability at most
$\varepsilon = 2^{-t}$. Hence
$\E_\xi \abs{(X_\xi + Y_\xi)_z} \ge (1 - \varepsilon)(s_z + g - 1)$, which is at least $s_z - 1$
as soon as $\varepsilon s_z \le (1 - \varepsilon) g$. This holds because
$\alpha s_z \le 4\min(m, n)/d$ and $\varepsilon \le \alpha\gamma_0/128$, which give
$\varepsilon s_z \le \theta/32 \le (1 - \varepsilon)\theta/16$.

Summing \eqref{eq:class} over $z$, and using $\abs{X} \le d \abs{X_{r_0}}$,
\[
  \E_\xi \abs{X_\xi + Y_\xi}
  \ge \sum_z \bigl(\abs{X_{r_0}} + \abs{Y_{z - r_0}} - 1\bigr)
  = d \abs{X_{r_0}} + \abs{Y} - d
  \ge \abs{X} + \abs{Y} - d . \qedhere
\]
\end{proof}

The proof of the paper uses \eqref{eq:pair} in the same two cases, with a random member of
$\mathcal F_{z - r_0}$ or $\mathcal F_{z - \xi_i}$ in place of $G(x, y)$, and takes the expectation
over these members as well. Since $x$ is $r_0$ in the first case and $\xi_i$ in the second, a set
$G(x, y)$ is needed only for $x \in S(\xi)$, and this is what the union over $S(\xi)$ provides.

\subsection*{Completion}
Fix $\xi$ with $\abs{X_\xi + Y_\xi} \ge \abs{X} + \abs{Y} - d$. As in the paper, adding at most
$d - 1$ further elements of $X$ to $X_\xi$ and of $Y$ to $Y_\xi$ raises the size of the sumset to
$\abs{X} + \abs{Y} - 1$, and enlarging the two sets to $c_1$ and $c_2$ elements keeps this bound.
This requires
\[
  (t + 1)\Bigl(\frac{m}{d} + 2\Bigr) + d - 1 \le c_1 ,
  \qquad
  d\,(t + 1)\,c_0 + d - 1 \le c_2 .
\]
With $K = \max\bigl((t + 1)c_0, \ceil{128/\gamma_0}\bigr)$, $d = \floor{c_2/(K + 1)}$ and
\[
  c = 16\,(t + 1)(K + 1),
\]
the hypothesis $c_1 c_2 \ge c \max(\abs{X}, \abs{Y})$ gives both inequalities and
$\theta \ge 128$. With Theorem 5 and the paper's sample of $Y$, the same argument works with
$K = \max\bigl(c_0, \ceil{128/\gamma_0}\bigr)$; replacing $c_0$ by $(t + 1)c_0$ is the only
change in the constants.

\section{Summary}

\begin{itemize}
\item Theorem 4 is Corollary 24 of the small paper and is used unchanged in Theorems 1 and 6.
\item Theorem 5 is replaced by Theorem~\ref{thm:weak5}, in which the bounded subset of $B$ may
  depend on $A$. It follows from the form of Theorem 31 proved in the small paper by reduction
  modulo a large prime, with the Cauchy--Davenport inequality in $\Z$ when $B$ is small.
\item In the proof of Theorem 8 the sample of $Y$ is the union of the sets $G(x, y)$ over the
  classes $x$ of the sample of $X$ and the dense classes $y$ of $Y$. The proof of Lemma F\&G is
  otherwise unchanged; the sample of $Y$ is $t + 1$ times larger, and the constant of Theorem 8
  is $16(t + 1)(K + 1)$ with $K = \max\bigl((t + 1)c_0, \ceil{128/\gamma_0}\bigr)$.
\item Theorem 1 follows from Theorems 4, 6 and 8 as in the paper.
\end{itemize}

\end{document}
